\documentclass[tikz,border=0pt]{standalone}
\usepackage[UTF8,fontset=none]{ctex}
\setCJKmainfont[BoldFont={Hiragino Sans GB W6}]{Hiragino Sans GB W3}
\setCJKsansfont[BoldFont={Hiragino Sans GB W6}]{Hiragino Sans GB W3}
\usetikzlibrary{arrows.meta,calc}
\definecolor{Canvas}{HTML}{FBFCFD}
\definecolor{Ink}{HTML}{172B2F}
\definecolor{Slate}{HTML}{52666A}
\definecolor{Brand}{HTML}{25636A}
\definecolor{col0fill}{HTML}{EAF6F7}
\definecolor{col0stroke}{HTML}{25636A}
\definecolor{col0text}{HTML}{194E54}
\definecolor{col1fill}{HTML}{FFF3E8}
\definecolor{col1stroke}{HTML}{D97745}
\definecolor{col1text}{HTML}{9A4B26}
\definecolor{col2fill}{HTML}{EDF7EF}
\definecolor{col2stroke}{HTML}{4F8A5B}
\definecolor{col2text}{HTML}{315D39}
\definecolor{col3fill}{HTML}{F5EEFA}
\definecolor{col3stroke}{HTML}{8E5AA7}
\definecolor{col3text}{HTML}{613774}
\begin{document}
\begin{tikzpicture}[
  x=1cm,
  y=1cm,
  every node/.style={font=\sffamily},
  card/.style={rounded corners=7pt, minimum height=2.15cm, inner sep=6pt, align=center, line width=0.8pt},
  flow/.style={-{Stealth[length=5pt,width=4pt]}, line width=1.2pt, draw=Brand!75, rounded corners=5pt}
]
\path[use as bounding box] (0,0) rectangle (25,10);
\fill[Canvas] (0,0) rectangle (25,10);
\fill[Brand] (0,9.78) rectangle (25,10);
\node[anchor=west, text=Ink] at (1,9.20) {\fontsize{18}{21}\selectfont\bfseries 第 10 章 · 评测与研究设计};
\node[anchor=west, text=Slate] at (1,8.55) {\fontsize{9.2}{11}\selectfont 从科学指标与系统基准，到失效归因、可证伪改进与新模型设计};
\draw[Brand!18, line width=0.7pt] (1,8.13) -- (24,8.13);
\node[card, fill=col0fill, draw=col0stroke, text width=5.12cm] (s1) at (3.748,4.75) {%
  {\fontsize{8.3}{10}\selectfont\bfseries\color{col0text} 10.1\quad 怎样评价世界模型}\par\vspace{2pt}
  {\fontsize{6.6}{8.299999999999999}\selectfont\color{Slate} 单步误差、长程漂移与任务效用}
};
\node[card, fill=col1fill, draw=col1stroke, text width=5.12cm] (s2) at (9.582,4.75) {%
  {\fontsize{8.3}{10}\selectfont\bfseries\color{col1text} 10.2\quad 世界模型的系统评测}\par\vspace{2pt}
  {\fontsize{6.6}{8.299999999999999}\selectfont\color{Slate} 开环推演、闭环控制与 OOD 鲁棒性}
};
\node[card, fill=col2fill, draw=col2stroke, text width=5.12cm] (s3) at (15.418,4.75) {%
  {\fontsize{8.3}{10}\selectfont\bfseries\color{col2text} 10.3\quad 失效分析与模型改进}\par\vspace{2pt}
  {\fontsize{6.6}{8.299999999999999}\selectfont\color{Slate} 幻觉、复合误差与协变量偏移}
};
\node[card, fill=col3fill, draw=col3stroke, text width=5.12cm] (s4) at (21.252,4.75) {%
  {\fontsize{8.3}{10}\selectfont\bfseries\color{col3text} 10.4\quad 设计新的世界模型}\par\vspace{2pt}
  {\fontsize{6.6}{8.299999999999999}\selectfont\color{Slate} 可证伪假设与多模态新范式}
};
\draw[flow] (s1.east) -- (s2.west);
\draw[flow] (s2.east) -- (s3.west);
\draw[flow] (s3.east) -- (s4.west);
\node[anchor=west, rounded corners=4pt, fill=Brand!8, text=Brand, inner xsep=7pt, inner ysep=4pt] at (1,0.62) {\fontsize{7.2}{9}\selectfont\bfseries 学习路径};
\node[anchor=west, text=Slate] at (3.25,0.62) {\fontsize{7.2}{9}\selectfont 概念建模 → 核心方法 → 工程实现 → 系统验证};
\node[anchor=east, text=Slate!72] at (24,0.62) {\fontsize{6.6}{8}\selectfont HANDS-ON WORLD MODELS};
\end{tikzpicture}
\end{document}
